\documentclass[10pt,twocolumn]{article}
\usepackage[T1]{fontenc}
\usepackage[utf8]{inputenc}
\usepackage{lmodern}
\usepackage[margin=1.8cm,top=2.4cm,bottom=2cm,columnsep=0.6cm]{geometry}
\usepackage{fancyhdr}
\usepackage{titlesec}
\usepackage{booktabs}
\usepackage{amsmath,amssymb,amsfonts}
\usepackage{microtype}
\usepackage{xcolor}
\usepackage{mdframed}
\usepackage{parskip}
\usepackage{enumitem}
\usepackage{hyperref}

\definecolor{rulered}{RGB}{180,30,30}
\definecolor{boxgray}{RGB}{245,245,245}
\definecolor{keywordgray}{RGB}{100,100,100}

\pagestyle{fancy}
\fancyhf{}
\fancyhead[L]{\textsc{\small Claude Mag}}
\fancyhead[C]{\small\textit{Machine Learning Research Digest}}
\fancyhead[R]{\small 2026-08-10}
\fancyfoot[C]{\small\thepage}
\renewcommand{\headrulewidth}{0.8pt}

\titleformat{\section}{\normalsize\bfseries\color{rulered}}{}{0pt}{}%
  [\vspace{-4pt}\color{rulered}\rule{\columnwidth}{0.4pt}]
\titlespacing{\section}{0pt}{8pt}{4pt}

\hypersetup{colorlinks=true,urlcolor=rulered,linkcolor=black}

\begin{document}
\setlength{\parindent}{0pt}
\setlength{\parskip}{4pt}

{\small\color{keywordgray}\textbf{Keywords:}
  3D vision-language model \textbullet{}
  spatial understanding \textbullet{}
  3D rotary positional embeddings \textbullet{}
  multi-view fusion}\par\vspace{4pt}

{\LARGE\bfseries\raggedright
  Qwen-3D: A Generalist 3D Vision-Language Model
  for Spatial Understanding}\par\vspace{4pt}

{\large\itshape\raggedright
  Carnegie Mellon Researchers Unify 2D and 3D Grounding
  by Computing Attention in World Space Rather Than Image Space}\par\vspace{6pt}

{\small\textbf{By} Lucy Lin, Ayush Jain, Yifan Liu,
  Katerina Fragkiadaki (Carnegie Mellon University)}\par\vspace{2pt}

{\small\color{keywordgray}arXiv:2608.02980 \textbullet{} ECCV 2026}
\par\vspace{2pt}\noindent{\color{rulered}\rule{\columnwidth}{0.6pt}}\vspace{6pt}

A new generalist vision-language model from Carnegie Mellon University achieves
state-of-the-art results on both 2D and 3D spatial understanding benchmarks by a
single architectural change: computing attention in 3D world coordinates rather than
per-frame pixel coordinates. Qwen-3D surpasses existing 3D LMMs and several large
proprietary 2D models on grounding tasks, and will appear at ECCV 2026.

\begin{mdframed}[backgroundcolor=boxgray,linecolor=rulered,linewidth=1pt,
    innerleftmargin=6pt,innerrightmargin=6pt,
    innertopmargin=6pt,innerbottommargin=6pt]
\textbf{\small AT A GLANCE}\par\vspace{4pt}
\begin{description}[leftmargin=0pt,labelsep=4pt,itemsep=2pt,parsep=0pt,
    font=\small\bfseries]
  \item[Problem:] Frame-centric tokenization ignores the shared 3D structure
    underlying multi-view observations, limiting VLMs on long videos and 3D scenes.
  \item[Why it matters:] Embodied AI and robotics need a single model that
    reasons about 3D space and natural language without separate specialist pipelines.
  \item[Method:] 3D Rotary Positional Embeddings index attention by world-space
    coordinates; a unified output head predicts language, 2D boxes, and 3D boxes.
  \item[Result:] State-of-the-art on ScanQA, EmbodiedScan, and other 3D benchmarks;
    competitive with large proprietary 2D models on MMBench and SEED-Bench.
  \item[Evidence:] Evaluated across six 3D understanding and five 2D vision-language
    benchmarks at ECCV 2026.
\end{description}
\end{mdframed}

\section{The Big Picture}

Large Multimodal Models have transformed image and short-video understanding, but
scaling to 3D scenes and long videos has remained elusive. Every RGB frame is
tokenized independently, discarding the geometric relationships that tie views of the
same scene together. A model watching a robot navigate a room from multiple cameras
sees a stream of disconnected image patches, with no built-in sense that the sofa
seen from Camera~1 is the same object partially visible in Camera~2.

Qwen-3D solves this by lifting attention from 2D pixel space into 3D world space.
Given depth maps and camera poses alongside RGB images, the model maps every visual
token to a 3D world coordinate and uses that coordinate as the token's positional
encoding. Tokens that correspond to the same 3D point---even from different cameras
or timesteps---receive encodings that reflect their world-space proximity, enabling
the model to naturally associate them during attention without any explicit
correspondence algorithm.

\section{Why Existing Approaches Fall Short}

Prior 3D LMMs (e.g.\ 3D-VisTA, EmbodiedScan) added 3D-awareness through auxiliary
encoders or separate perception heads, but kept the core transformer operating in
2D image space. This created a mismatch: the language backbone reasons over 2D
feature sequences while 3D grounding heads operate over a separate 3D feature space,
requiring careful multi-task balancing and making joint 2D/3D generalization difficult.

Specialist 3D perception models (Mask3D, OpenMask3D) achieve strong grounding results
but are not language models and cannot engage in open-ended dialogue or instruction
following. Generalist VLMs (GPT-4V, Gemini Pro Vision) handle language fluently but
treat 3D scenes as collections of independent image crops, losing spatial coherence.
Qwen-3D is the first model to close all three gaps in a single architecture.

\section{Core Mechanism}

\textbf{3D Token Lifting.}
For each pixel $(u, v)$ in an input frame, depth $d$ and camera intrinsics
$K$ and extrinsics $[R | t]$ yield the world coordinate:
\[
\mathbf{p} = R^{-1}\!\left(d\,K^{-1}\begin{bmatrix}u\\v\\1\end{bmatrix} - t\right).
\]
Each visual token (a patch of pixels) is assigned the 3D centroid of its patch.

\textbf{3D Rotary Positional Embeddings (3D RoPE).}
Standard RoPE encodes 1D or 2D positions by rotating query and key vectors in
frequency-indexed subspaces. Qwen-3D extends this to 3D: the world coordinates
$(X, Y, Z)$ are used to construct rotation matrices for three independent groups
of attention head dimensions, one per spatial axis. Two tokens close in 3D world
space receive nearly identical positional rotations, producing high raw dot-product
similarity before any content-based weighting.

\textbf{Unified Output Head.}
A single transformer decoder produces three output types from the same representation:
free-form text, 2D bounding boxes (normalized image coordinates), and 3D bounding
boxes (centre, size, heading angle in world coordinates). Task type is controlled by
instruction prompts, enabling zero-shot switching between modalities.

\section{Technical Formulation}

Let $\mathbf{F} = \{(\mathbf{p}_i, \mathbf{v}_i)\}_{i=1}^N$ be the set of $N$ visual
tokens with 3D world positions $\mathbf{p}_i \in \mathbb{R}^3$ and feature vectors
$\mathbf{v}_i \in \mathbb{R}^d$. Self-attention with 3D RoPE applies:
\[
\mathrm{Attn}(Q, K, V)_{ij}
= \frac{\bigl(R(\mathbf{p}_i)\,q_i\bigr)^{\!\top}
         \bigl(R(\mathbf{p}_j)\,k_j\bigr)}{\sqrt{d_h}},
\]
where $R(\mathbf{p}) \in \mathbb{R}^{d_h \times d_h}$ is the block-diagonal rotation
matrix constructed from the three spatial frequencies of $\mathbf{p}$.

The training objective combines cross-entropy for text tokens ($\mathcal{L}_\text{LM}$),
smooth-$\ell_1$ for 2D box coordinates ($\mathcal{L}_\text{2D}$), and smooth-$\ell_1$
plus IoU3D loss for 3D box prediction ($\mathcal{L}_\text{3D}$):
\[
\mathcal{L} = \mathcal{L}_\text{LM} + \lambda_1\,\mathcal{L}_\text{2D}
              + \lambda_2\,\mathcal{L}_\text{3D}.
\]

\section{Learning or Inference Procedure}

\textbf{Data.} Training uses a mixture of (i) standard 2D VQA and instruction-following
datasets, (ii) multi-view indoor scene datasets with 3D annotations (ScanRefer, Nr3D),
and (iii) egocentric video datasets with depth and camera poses.

\textbf{Fine-tuning.} Starting from a Qwen-VL checkpoint, the model is fine-tuned
end-to-end with the multi-task objective above. The 3D RoPE replaces the existing 2D
RoPE positional encoding; no architectural changes to the attention mechanism itself
are required.

\textbf{Inference.} At test time, the model accepts any combination of RGB images,
depth maps, and camera poses. Depth-free inference falls back to 2D RoPE, preserving
2D benchmark performance even when 3D information is unavailable.

\section{What the Guarantee Says}

3D RoPE preserves the rotation equivariance of standard RoPE: if all tokens are
rigidly rotated in world space (equivalent to a global camera pose change), the
relative positional encoding between any pair of tokens is unchanged, making the model
invariant to global reference frame choice. The unified training objective ensures that
improving 3D grounding does not degrade language modeling performance---empirically
confirmed by stable scores on 2D benchmarks throughout training.

\section{Experimental Findings}

\begin{itemize}[itemsep=2pt,parsep=0pt]
  \item \textbf{ScanQA:} new state of the art, surpassing previous best 3D LMM by
    a significant margin on both EM and CIDEr metrics.
  \item \textbf{EmbodiedScan:} top result on 3D visual grounding (Acc@0.25 and Acc@0.5),
    outperforming specialist 3D perception pipelines.
  \item \textbf{MMBench / SEED-Bench:} competitive with GPT-4V and Gemini Pro Vision
    on standard 2D multimodal benchmarks, demonstrating no 3D/2D trade-off.
  \item \textbf{Multi-view video understanding:} large improvements on tasks requiring
    cross-view association (e.g.\ object re-identification across camera angles).
\end{itemize}

\section{Ablations and Interpretation}

\textbf{3D RoPE vs.\ 2D RoPE.} Removing 3D positional information and reverting to
per-frame 2D RoPE drops 3D grounding accuracy by over 15 points on ScanRefer,
confirming that world-space position is the key ingredient rather than multi-view data
alone.

\textbf{Depth quality sensitivity.} Using predicted depth (from a monocular estimator)
instead of ground-truth depth degrades 3D grounding by 4--6 points but still
outperforms prior 3D LMMs that use GT depth, suggesting the 3D positional signal is
robust to moderate depth noise.

\textbf{Joint vs.\ separate 2D/3D training.} Training only on 3D data degrades 2D
benchmark performance by 8 points; joint training recovers it fully, confirming that
the two modalities are complementary rather than competing.

\section*{Reference}

Lucy Lin, Ayush Jain, Yifan Liu, Katerina Fragkiadaki.
\textbf{Qwen-3D: A Generalist 3D Vision-Language Model for Spatial Understanding}.
arXiv:2608.02980, ECCV 2026.
\url{https://arxiv.org/abs/2608.02980}

\end{document}
